\chapter{The Results, and Why They Came Out That Way}
\label{ch:results}

This chapter takes each result in turn and answers three questions about it:
\textbf{what} was measured, \textbf{why} it turned out that way, and \textbf{how
much confidence} it deserves. Where a result replaced an earlier one, the earlier
one is shown too --- the direction in which a number moves when it is checked is
itself information.

\section{How performance is measured, and why not by success rate}

\begin{center}
\begin{tabular}{@{} L{4.0cm} L{10.0cm} @{}}
\toprule
\textbf{Metric} & \textbf{Definition and what it tells you} \\
\midrule
Final error & Hand-to-target distance at the last step. \textbf{The headline number:} it measures reaching \emph{and} holding. \\
Closest approach & Smallest distance at any instant. Separates ``can it get there'' from ``can it stay''. \\
Ends within 5 / 10\,cm & Graded proximity. Stable where a strict threshold is not. \\
Drift & Closest approach minus final error --- how much is given back after arriving. \\
Path length & Distance travelled per \SI{0.80}{\meter} of required travel. Exposes sweeping and overshoot. \\
Energy & Summed squared activation. Distinguishes a passive policy from a straining one. \\
Success rate & Within \SI{2}{\centi\meter} \emph{and} settled. \textbf{Reported, never headlined.} \\
\bottomrule
\end{tabular}
\end{center}

\begin{why}
\textbf{Why success rate is not the headline.} Across five runs of one identical
configuration it read 0\,\%, 4\,\%, 6\,\%, 8\,\% and 20\,\%. A measure that
varies twentyfold under identical settings is reporting sampling noise, not
performance.

The reason is structural: 4\,\% of fifty episodes is \emph{two episodes}. Any
binary threshold near the edge of a system's capability behaves this way. Final
error averages over all fifty and is stable to 4\,\% across the same five runs.
This is why the graded proximity measures exist --- they retain the
interpretability of a threshold without its brittleness.
\end{why}

\section{The self-destruction exploit}
\label{sec:defect3}

\textbf{What was measured.} Fifty evaluation episodes of a policy that appeared,
from its reward curve, to be learning.\ \textbf{Result: 100\,\% ended in a
deliberate numerical blow-up at step 7 of 100.} Not one episode ran to completion.

\begin{why}
\textbf{Why it happened.} Every per-step reward term was zero or negative, and a
blow-up terminated the episode with no penalty attached. In this framework a
terminated episode accrues nothing further --- its remaining value is exactly
zero --- while surviving all 100 steps accumulated about $-20$.

\textbf{Quitting was worth 20 more than participating.} The agent found the
cheapest way to quit: drive the muscles into a configuration the integrator
cannot handle.

The algorithm was not broken. It maximised the objective it was given, precisely.
The objective was wrong.
\end{why}

\textbf{The fix, and why it is the shape it is.} A survival term $w_4 = 4.5$ paid
\emph{every step}, chosen so the per-step reward is strictly positive at any
reachable error (costs are bounded by $2.7 + 1.0 + 0.5 = 4.2$), plus an explicit
blow-up penalty $w_6 = 10$. Living now strictly beats quitting from every state
--- not on average, but pointwise, which is the property that makes the fix
robust rather than lucky.

\textbf{Confidence: certain.} The failure was directly observed on 50 of 50
episodes and the blow-up rate is exactly zero across all five replication seeds
afterwards.

\section{The discount factor is the whole story}
\label{sec:gammaresult}

This is the most important result in the project, and it was arrived at by
first being wrong.

\subsection{The wrong answer, held for weeks}
\label{sec:entropyrefuted}

A sixteen-trial Bayesian search over five hyperparameters produced a
configuration that cut final error from \SI{0.261}{} to \SI{0.161}{\meter}. All
five best trials selected a target entropy near $-19.6$ against the library
default of $-9$.

\begin{why}
\textbf{Why that looked conclusive.} The mechanism fit the symptom exactly. SAC
pays its policy to stay random; residual randomness is precisely what stops an
end effector settling; the default scales with the \emph{number of muscles}
rather than with anything about the task; and every good trial independently
agreed on roughly half the default randomness.

Four independent lines of evidence, all pointing one way. It was written up as a
contribution of the thesis.
\end{why}

\begin{trap}
\textbf{It was wrong.} A two-sided one-at-a-time ablation, three seeds per arm:

\begin{center}
\begin{tabular}{@{} L{7.6cm} C{3.2cm} C{2.4cm} @{}}
\toprule
\textbf{Configuration} & \textbf{Final error (m)} & \textbf{Gap closed} \\
\midrule
All defaults \emph{(reference)} & $0.2613 \pm 0.0203$ & 0\,\% \\
Defaults $+$ tuned entropy only & $0.2688 \pm 0.0090$ & $-8$\,\% \\
Tuned $-$ entropy reverted & $0.1552 \pm 0.0181$ & $+106$\,\% \\
All tuned \emph{(reference)} & $0.1610 \pm 0.0084$ & 100\,\% \\
\bottomrule
\end{tabular}
\end{center}

Adding the tuned entropy to defaults recovers \emph{nothing}. Removing it from
the tuned configuration costs \emph{nothing}. Both directions agree, so this is
not a marginal call.

\textbf{Why the evidence misled.} A TPE sampler concentrates its proposals where
performance is good. Parameters therefore cluster in the best trials whether or
not they are causally responsible. Reading causation from that clustering is a
plain error --- and it is the error this project made.

\textbf{A hyperparameter search identifies a configuration, never a cause.} Only
a deliberate ablation separates the two.
\end{trap}

\subsection{The right answer}

The same procedure applied to all five parameters:

\begin{center}
\begin{tabular}{@{} L{6.8cm} C{3.4cm} C{2.4cm} @{}}
\toprule
\textbf{Defaults, plus one tuned parameter} & \textbf{Final error (m)} & \textbf{Gap closed} \\
\midrule
(nothing changed) & $0.2613 \pm 0.0203$ & 0\,\% \\
target entropy only & $0.2688 \pm 0.0090$ & $-8$\,\% \\
learning rate only & $0.2642 \pm 0.0026$ & $-3$\,\% \\
$\tau$ only & $0.2660 \pm 0.0117$ & $-5$\,\% \\
batch size only & $0.2658 \pm 0.0096$ & $-5$\,\% \\
\textbf{discount factor only} & $\mathbf{0.1553 \pm 0.0061}$ & $\mathbf{+106}$\,\% \\
\midrule
All tuned \emph{(reference)} & $0.1610 \pm 0.0084$ & 100\,\% \\
\bottomrule
\end{tabular}
\end{center}

\textbf{Five parameters tested individually. Four do nothing. One explains
everything.}

\begin{plain}
The discount factor decides how far into the future the agent bothers to look ---
about $1/(1-\gamma)$ steps ahead.

\begin{center}
\begin{tabular}{@{} L{5.4cm} C{2.4cm} C{5.0cm} @{}}
\toprule
& $\gamma$ & \textbf{effective horizon} \\
\midrule
Library default & $0.99$ & 100 steps --- the whole episode \\
Selected by the search & $0.957$ & 23 steps $=$ \SI{0.46}{\second} \\
\bottomrule
\end{tabular}
\end{center}

Now compare that with what the arm actually does. Measured from the trajectory,
the hand arrives at the target by \textbf{step 25} --- \SI{0.50}{\second}.

The tuned value matches the duration of the movement almost exactly. The default
was four times too long.
\end{plain}

\begin{why}
\textbf{Why a too-long horizon hurts, mechanically.} At $\gamma = 0.99$ the value
of an early step includes seventy-odd later steps of holding reward. That is a
high-variance quantity to estimate. A critic fed high-variance targets produces
unreliable estimates; an actor trained against an unreliable critic cannot refine
its endpoint. The visible symptom is an evaluation curve that oscillates and
never sharpens --- exactly what every untuned run in this project showed.

Shortening the horizon to the length of the reach makes credit assignment local
and the value estimate tractable.

\textbf{The transferable statement,} which is the part worth carrying to another
project: \emph{match the discount factor to the timescale of the behaviour, not
to a library default. Where the default horizon greatly exceeds the task
duration, the resulting variance presents as an inability to refine --- a symptom
easily and wrongly blamed on insufficient training, on reward shaping, or on
exploration.} This project spent considerable effort on all three before testing
$\gamma$.
\end{why}

\textbf{Confidence: high.} Single parameter changed, three seeds, full effect
recovered in both directions, four alternatives excluded by the identical
procedure, and a mechanism that predicts the observed timescale quantitatively.

\section{What the tuned policy actually does}

\begin{center}
\begin{tabular}{@{} L{5.6cm} C{3.0cm} C{3.0cm} @{}}
\toprule
\textbf{Metric (300k steps, 50 episodes)} & \textbf{Default} & \textbf{Tuned} \\
\midrule
Final error (m) & 0.204 & \textbf{0.106} \\
Closest approach (m) & 0.132 & \textbf{0.074} \\
Drift (m) & 0.072 & \textbf{0.035} \\
Ends within \SI{10}{\centi\meter} & 18\,\% & \textbf{58\,\%} \\
Energy & 774\,667 & \textbf{276\,804} \\
Path length per \SI{0.80}{\meter} reach & \SI{3.55}{\meter} & \textbf{\SI{1.64}{\meter}} \\
\bottomrule
\end{tabular}
\end{center}

\begin{plain}
Read the path-length row. To cover \SI{0.80}{\meter} of required travel, the
untuned policy moved its hand \SI{3.55}{\meter} --- it swept around, overshot,
and came back. The tuned policy moved \SI{1.64}{\meter}, with net displacement
$0.99\times$ the distance required.

It stopped wandering and started going there.
\end{plain}

Intra-episode, averaged over thirty episodes: the hand starts \SI{0.795}{\meter}
away, is at \SI{0.198}{} by step 10, \SI{0.119}{} by step 25, and then holds ---
\SI{0.110}{}, \SI{0.112}{}, \SI{0.113}{\meter} at steps 50, 75, 99. \textbf{That
flat tail is the whole point:} converged and stationary, not still drifting.

\section{Comparing the four algorithms, and getting it wrong first}
\label{sec:fairness}

\textbf{What was measured first.} Five configurations, three seeds, 100k steps.
Ordering: SAC $<$ TD3 $<$ DDPG $<$ PPO. DDPG was reported as clearly the weakest
off-policy method, explained by its lack of twin critics --- the textbook account.

\begin{trap}
\textbf{That comparison was invalid, and by its own later findings.} SAC ran with
$\gamma = 0.957$; the other three ran at the library default $0.99$. Once
$\gamma$ was shown to be the decisive parameter, the benchmark stood revealed as
\emph{one correctly configured entrant against three handicapped ones}.

Retrained with $\gamma = 0.957$ for all four:

\begin{center}
\begin{tabular}{@{} L{2.8cm} C{3.0cm} C{3.4cm} C{2.4cm} @{}}
\toprule
& \textbf{Before ($\gamma{=}0.99$)} & \textbf{After ($\gamma{=}0.957$)} & \textbf{Change} \\
\midrule
SAC  & $0.2613 \pm 0.0203$ & $\mathbf{0.1909 \pm 0.0341}$ & $-27$\,\% \\
DDPG & $0.4161 \pm 0.0189$ & $\mathbf{0.2510 \pm 0.0383}$ & $\mathbf{-40}$\,\% \\
TD3  & $0.2786 \pm 0.0110$ & $0.2579 \pm 0.0348$ & $-7$\,\% \\
PPO  & $0.4231 \pm 0.0351$ & $0.4661 \pm 0.0401$ & $+10$\,\% \\
\bottomrule
\end{tabular}
\end{center}

DDPG improves 40\,\% and overtakes TD3. The ranking changes. \textbf{The
twin-critic explanation was constructed after the fact around a result that has
not survived} --- DDPG was not failing for want of twin critics; it was failing
because its horizon was wrong, and it suffered more from that than the others.
Corrected, it is indistinguishable from TD3 and no ordering between them is
supported.
\end{trap}

\begin{why}
\textbf{Why PPO got worse}, uniquely. PPO estimates advantages through
generalised advantage estimation, where $\gamma$ and $\lambda$ \emph{jointly} set
the bias--variance trade-off. Lowering $\gamma$ while leaving $\lambda = 0.95$
shortens the horizon on one side of a pair only. The correct adjustment for PPO
touches both, and was not attempted --- so PPO's last place is again about
configuration, not about the algorithm.

This is what makes fair comparison of RL algorithms genuinely hard: there is no
single setting that can be held equal across all four, because the parameters do
not mean the same thing in each.
\end{why}

\textbf{The conclusion that survives.} The gap between the best and second-best
off-policy method is \SI{0.060}{\meter}. The improvement obtained in DDPG by
correcting $\gamma$ alone was \SI{0.165}{\meter} --- roughly \textbf{three times
the gap between algorithms}.

\begin{center}
\begin{tabular}{@{}p{3pt}@{\hspace{10pt}}p{0.93\linewidth}@{}}
\cellcolor{navy} & On this task, choosing the discount factor correctly mattered
about three times more than choosing between SAC, TD3 and DDPG.
\end{tabular}
\end{center}

\textbf{Confidence: moderate.} Three seeds, no significance testing (three seeds
cannot support a signed-rank test with useful power, so none is claimed), a
100k-step early-training snapshot, and only $\gamma$ corrected --- a \emph{fair}
comparison, not a \emph{tuned} one. A tuned comparison needs a separate search
per algorithm, which the compute budget did not allow.

\section{The central result: do the muscle forces agree?}

This is the thesis question. Everything above concerns whether the arm reaches;
this concerns whether it reaches \emph{the same way}.

\subsection{How the comparison is made fair}

At every timestep of a rollout, take the joint torque $\tau$ the policy actually
produced, hand it to static optimisation, and ask what forces \emph{it} would
have chosen for that same torque.

\begin{why}
\textbf{Why $\tau$ is taken from the policy rather than computed independently.}
Two reasons, and both are about eliminating excuses.

It is \emph{exact} --- \texttt{qfrc\_actuator} is the generalised force the
muscles genuinely produced, so no numerical differentiation enters. And it is
\emph{guaranteed feasible} --- the policy's own force vector satisfies the
constraint by construction, so static optimisation can never fail for want of a
solution.

Both methods are then posed an identical problem on an identical trajectory. Any
difference is purely a difference of load-sharing strategy, which is exactly what
the thesis asks about.

\textbf{The cost of this choice, stated plainly:} the classical method never
chooses a trajectory of its own. So this validates load sharing \emph{along a
given path}, not the path. A controller could move in a way no human would and
still share load classically. Trajectory realism is not assessed anywhere in this
work.
\end{why}

\textbf{The identity was verified, not assumed.} Over 2000 timesteps, the mean
residual $\|R^\top f - \tau\|$ is \SI{1.46e-15}{\newton\meter} for the policy's
forces and \SI{1.68e-15}{} for the classical solution, with 0 optimisation
failures. Both at machine precision. A sign error or a moment-arm extraction
error would have invalidated the whole chapter, and would have shown up here.

\subsection{Result one: the pattern agrees}

\begin{center}
\begin{tabular}{@{} L{8.2cm} R{4.0cm} @{}}
\toprule
Pearson correlation, all muscles & $\mathbf{0.81 \pm 0.04}$ \\
Pattern cosine similarity & $0.813 \pm 0.047$ \\
\quad null baseline & $0.372 \pm 0.017$ \\
\quad \textbf{excess over null} & $\mathbf{+0.441}$ \\
\bottomrule
\end{tabular}
\end{center}

\begin{trap}
\textbf{The cosine number cannot be quoted on its own, and this nearly went
wrong.} Muscle forces are non-negative by construction, so two \emph{completely
unrelated} force vectors already agree strongly --- all their components point
into the same corner of the space.

The correct null: at each timestep, permute which muscle each of the policy's own
force magnitudes belongs to. This destroys the correspondence while preserving
that instant's magnitude distribution exactly.

That null averages \textbf{0.373}. So of an apparent 0.813, roughly the first
0.37 is a floor imposed by arithmetic, not evidence of anything.

\textbf{A second-order error inside the first.} The null was initially computed on
\emph{time-averaged} force vectors and gave 0.767 --- which would have made the
result look like nothing at all. That was also wrong: averaging over time removes
the very variation the statistic is meant to test. Computed per timestep, as it
must be, it is 0.373.

And because the null's 95th percentile for \emph{individual} timesteps is 0.867,
no single timestep's cosine means anything. Only the aggregate does.

\textbf{The Pearson correlation is mean-centred, carries no such floor, and is
the figure to quote when one number is required.}
\end{trap}

\subsection{Result two: the economy does not agree}

Producing identical joint torques, the learned controller expends
$\mathbf{1.91 \pm 0.15}$ times the effort of the minimum-effort solution.

\textbf{Where the excess sits} --- it is not spread evenly:

\begin{center}
\begin{tabular}{@{} L{3.6cm} R{2.2cm} R{2.4cm} R{1.6cm} R{1.8cm} @{}}
\toprule
\textbf{Muscle} & \textbf{Policy (N)} & \textbf{Classical (N)} & \textbf{$r$} & \textbf{NRMSE} \\
\midrule
deltoid anterior  &  63.7 &  59.2 & 0.959 &  2.5\,\% \\
deltoid medial    & 145.3 & 112.3 & 0.954 &  6.1\,\% \\
deltoid posterior &  35.5 &  38.7 & 0.859 & 10.0\,\% \\
triceps long      & 180.4 & 134.0 & 0.910 & 11.2\,\% \\
biceps long       &  52.1 &  41.6 & 0.893 &  7.5\,\% \\
brachialis        & 113.9 &  87.6 & 0.871 &  8.6\,\% \\
biceps short      &  90.7 &  45.7 & 0.764 & 18.1\,\% \\
\textbf{triceps lateral} &  79.1 &  25.8 & \textbf{0.416} & 16.5\,\% \\
\textbf{triceps medial}  &  64.3 &  14.2 & \textbf{0.416} & 15.3\,\% \\
\bottomrule
\end{tabular}
\end{center}

\begin{plain}
The three shoulder muscles agree closely --- one to within 2.5\,\%. The problem is
in two elbow extensors, where the policy uses three to four and a half times what
is prescribed.

Those muscles \emph{straighten} the elbow while the biceps group \emph{bends} it.
Driving both at once produces largely cancelling torques. The elbow ends up at
roughly the right angle, but both groups have spent large forces to achieve what
one alone would have done.

This is \textbf{co-contraction}. Humans do it deliberately when a joint needs to
be stiff --- carrying a full cup, or steadying a hand. Doing it throughout an
unloaded reach is simply waste, and it is not what the minimum-effort criterion
selects.
\end{plain}

\textbf{Independently corroborated.} The Falconer--Winter co-contraction index
over the elbow flexor/extensor groups reads 0.564 --- elevated, and the only
behavioural metric that did \emph{not} improve under tuning (0.501 before). Two
methodologically unrelated analyses, one a comparison against classical
optimisation and one a standard EMG-derived index, localise the same defect to
the same muscle group. That agreement is stronger evidence than either alone.

\subsection{Result three: the task is solved consistently, the muscle solution is not}

Five identical runs differing only in random seed:

\begin{center}
\begin{tabular}{@{} C{1.6cm} C{2.2cm} C{2.2cm} C{2.0cm} C{1.8cm} C{1.8cm} @{}}
\toprule
\textbf{Seed} & \textbf{Final error} & \textbf{Closest} & \textbf{Effort ratio} & \textbf{$r$} & \textbf{Success} \\
\midrule
0 & 0.1142 & 0.0766 & 1.71 & 0.865 & 0.04 \\
1 & 0.1086 & 0.0785 & 1.91 & 0.823 & 0.20 \\
2 & 0.1034 & 0.0757 & 2.12 & 0.759 & 0.00 \\
3 & 0.1066 & 0.0672 & 1.86 & 0.825 & 0.08 \\
4 & 0.1041 & 0.0691 & 1.97 & 0.767 & 0.06 \\
\midrule
\textbf{mean} & $0.1074$ & $0.0734$ & $1.914$ & $0.808$ & $0.076$ \\
\textbf{std}  & $0.0043$ & $0.0050$ & $0.150$ & $0.044$ & $0.075$ \\
\bottomrule
\end{tabular}
\end{center}

Two facts sit side by side in that table and their combination is the most
informative result in the thesis. \textbf{Final error varies by 4\,\% across the
five runs. Muscular energy varies by 77\,\%} (273{,}241 to 483{,}373).

\begin{plain}
Five controllers, trained identically apart from a random seed, all arrive at
essentially the same place --- using materially different sets of muscle forces.

That is the redundancy problem of Chapter~1 appearing directly in the results.
Nine muscles, four joints; many force distributions give the same movement; and
nothing in the reward chose among them. The learner settled somewhere different
each time.

It also tells you where the fix is. Huge variance in the muscular solution
alongside near-constancy in the task outcome is the signature of \textbf{an
objective that is undetermined in exactly the dimension you care about.}
\end{plain}

\begin{trap}
\textbf{The originally reported run was the best of the five.} Seed 0 --- the run
every earlier section analysed --- gave the \emph{lowest} effort ratio (1.71 vs a
mean of 1.914) and the \emph{highest} correlation (0.865 vs 0.808). On both
metrics carrying the argument, the single run first reported was the most
favourable of five.

Not a large effect, but a systematic one, and exactly the error that reporting a
single run invites. Without replication this thesis would have reported both an
agreement and an efficiency better than the method delivers.
\end{trap}

\section{Closing the gap: the effort weight}
\label{sec:sweepresult}

\textbf{The prediction, made before the experiment.} If the effort term is 28
times smaller than the precision bonus, it cannot be influencing the solution ---
so raising $w_2$ should reduce the discrepancy, at some cost in accuracy.

Run at $w_2 \in \{1, 5, 20\}$, everything else fixed, five seeds each, full
300k-step budget, and evaluated under the \emph{default} weights so the numbers
stay comparable with every other result:

\begin{center}
\begin{tabular}{@{} L{2.4cm} C{2.6cm} C{2.4cm} C{2.6cm} C{2.0cm} @{}}
\toprule
\textbf{Config.} & \textbf{Effort ratio} & \textbf{Pearson $r$} & \textbf{Final error} & \textbf{Energy} \\
\midrule
$w_2 = 1$  & $1.914 \pm 0.150$ & $0.808 \pm 0.044$ & $0.1074 \pm 0.0079$ & 327\,048 \\
$w_2 = 5$  & $1.406 \pm 0.092$ & $0.919 \pm 0.016$ & $0.1114 \pm 0.0053$ & 104\,598 \\
$\mathbf{w_2 = 20}$ & $\mathbf{1.304 \pm 0.025}$ & $\mathbf{0.943 \pm 0.008}$ & $0.1161 \pm 0.0108$ & \textbf{37\,076} \\
\midrule
Conventional & $1.263$ & $0.931$ & $0.165$ & 102\,876 \\
\bottomrule
\end{tabular}
\end{center}

The prediction held in the means. The per-seed ranges tell a more careful story,
and it is worth getting right.

\begin{center}
\begin{tabular}{@{} L{2.2cm} L{7.2cm} C{3.0cm} @{}}
\toprule
& \textbf{the five seeds} & \textbf{range} \\
\midrule
$w_2 = 1$  & 1.714, 1.856, 1.912, 1.974, 2.116 & 1.714--2.116 \\
$w_2 = 5$  & \textbf{1.336}, 1.351, 1.355, 1.429, 1.557 & 1.336--1.557 \\
$w_2 = 20$ & 1.278, 1.292, 1.293, 1.318, \textbf{1.345} & 1.278--\textbf{1.345} \\
\bottomrule
\end{tabular}
\end{center}

\begin{trap}
An earlier version of this document said the three distributions were
\emph{cleanly separated}. Checked against the per-seed numbers, that is true of
one step and false of the other.

$w_2 = 1 \rightarrow 5$ \textbf{is} clean: every seed of one lies below every seed
of the other, $1.557 < 1.714$.

$w_2 = 5 \rightarrow 20$ \textbf{is not}: the worst $w_2 = 20$ seed at
$\mathbf{1.345}$ sits just above the best $w_2 = 5$ seed at $\mathbf{1.336}$. The
ranges overlap by 0.009.

The means still move the right way, and four of the five $w_2 = 20$ seeds do
clear the entire $w_2 = 5$ range --- but one overlapping pair at $n = 5$ makes
that second step a \emph{probable} further improvement, not a demonstrated one.
The upper bound had also been mis-transcribed as 1.340.
\end{trap}

Per-muscle RMSE falls from 10.1\,\% to about 3\,\% of mean $F_{\max}$.

\textbf{The accuracy cost is real but modest}: 8\,\% worse final error for a
32\,\% reduction in wasted effort. Closest approach is unchanged
(0.0734, 0.0750, 0.0737), so what is lost is \emph{holding}, not \emph{reaching}.

\begin{why}
\textbf{The most informative row is the standard deviation.} It collapses:
$0.150 \rightarrow 0.092 \rightarrow 0.025$ on effort, and
$0.044 \rightarrow 0.016 \rightarrow 0.008$ on correlation.

That confirms the interpretation offered above. At $w_2 = 1$ the muscular
solution varied 77\,\% across seeds while task performance stayed constant,
because the objective did not constrain it. Weighting that dimension \emph{pins
the solution down}: at $w_2 = 20$ the load-sharing outcome is determined by the
objective rather than by the random seed.

This is a stronger form of evidence than the mean improvement alone. It shows the
mechanism, not just the effect.
\end{why}

\begin{center}
\begin{tabular}{@{}p{3pt}@{\hspace{10pt}}p{0.93\linewidth}@{}}
\cellcolor{navy} & \textbf{The headline answer changes.} It is not that DRL
reproduces coordination but not economy. It is that DRL reproduces
\textbf{both} --- provided the objective asks for both. The version originally
reported did not ask.
\end{tabular}
\end{center}

\subsection{Two predictions, one right and one wrong}

Both are recorded because the direction of the error is instructive.

\textbf{Right:} the prediction that raising $w_2$ would trade accuracy for
fidelity, made in advance, borne out in both directions.

\textbf{Wrong:} this author predicted that $w_2 = 20$ would \emph{regress} on
replication, reasoning that single-seed results in this project had twice proved
optimistic. It did the opposite --- seed 0's 1.34 was the \emph{worst} of the
five and the mean is 1.304. Selection bias explains a favourable first
measurement; it does not license assuming that every first measurement is
favourable.

\section{Against a conventional controller --- and a withdrawn claim}

Static optimisation is an \emph{idealised} reference: solved instant by instant,
with no actuator lag and no tracking. Comparing against it says how far the
policy is from a mathematical optimum, not how far it is from what conventional
control actually achieves --- and the thesis question is the latter.

A Cartesian impedance controller was therefore built:
$f_{\text{des}} = K_p(x^\star - x) - K_d\dot{x}$, mapped to joint torques by the
Jacobian transpose, with gravity and Coriolis terms compensated, and distributed
onto muscles by minimum-effort load sharing over the achievable force range. Both
gains were grid-searched, because a tuned policy against an untuned controller is
not a test.

\begin{trap}
\textbf{Two defects would have made it a strawman.} Both are on the
\emph{classical} side, and both were found and fixed before publishing anything.

\textbf{No gravity compensation.} Without the bias term the arm droops until
position error alone generates enough torque, leaving a large steady-state
offset. Measured, it reached only \SI{0.46}{\meter} --- \emph{worse than a random
policy}. Any textbook impedance controller includes this term.

\textbf{It assumed force is proportional to activation.} It is not: force is
$g(l,v)\,a + b(l,v)$, state-dependent. Converting a desired force to an
activation requires inverting the muscle model at the current length and velocity.
Probing at $a = 0$ and $a = 1$ recovers that affine relation exactly.

Reporting the comparison without these fixes would have overstated the case for
learning. This is the single easiest way to fake a machine-learning result, and
it is usually done by accident.
\end{trap}

\subsection{Why the conventional controller is not 1.00}

It solves minimum-effort load sharing explicitly at every timestep, so its effort
ratio ``should'' be exactly 1.00. It is \textbf{1.263}.

\begin{why}
The gap is not an error --- it is the price of closed-loop control. Activation
dynamics mean the commanded activation is not the realised one; the controller
acts at \SI{50}{\hertz} rather than continuously; and it tracks a moving state
rather than solving a static instant.

\textbf{The idealised 1.00 is not attainable by any controller operating under
these constraints}, and this reframes the central number. Measured against 1.00,
the learned policy appears to waste 91\,\% of its effort. Measured against 1.263 ---
the achievable floor, demonstrated by a controller that optimises effort
explicitly --- it uses about \textbf{51\,\% more than necessary}. The second figure
is the meaningful one.
\end{why}

\subsection{A claim made, then withdrawn}

\begin{trap}
An earlier draft observed that the learned policy at $w_2 = 20$ reached
$1.304 \pm 0.025$, \emph{below} the 1.33 then measured for the conventional
controller --- and cautioned that the 1.33 rested on a single gain pair from a
coarse twelve-point grid, so a better-tuned controller might overturn it.

A thirty-point search was run. It did. The coarse search had bottomed out against
its own lower boundary; the better gains are \emph{gentler} in both stiffness and
damping ($k_p = 20$, $k_d = 6$ against $k_p = 60$, $k_d = 25$), and give both
better accuracy (\SI{0.165}{} vs \SI{0.199}{\meter}) and better economy (1.263 vs
1.330).

\textbf{The claim that a learned policy distributes force more economically than
a controller optimising that quantity in closed form is therefore not supported,
and the paragraph advancing it was removed.}
\end{trap}

\textbf{The honest comparison, both sides measured properly:}

\begin{center}
\begin{tabular}{@{} L{5.6cm} C{3.4cm} C{3.4cm} @{}}
\toprule
& \textbf{Learned ($w_2 = 20$)} & \textbf{Conventional} \\
\midrule
Final error & $\mathbf{0.1072 \pm 0.0148}$ & \SI{0.165}{\meter} \\
Effort ratio & $1.304 \pm 0.025$ & $\mathbf{1.263}$ \\
Agreement with classical ($r$) & $\mathbf{0.943 \pm 0.008}$ & $0.931$ \\
\bottomrule
\end{tabular}
\end{center}

Neither method dominates. The learned policy reaches about 35\,\% more
accurately; the conventional controller is about 3\,\% more economical --- a
margin comparable to the learned policy's own seed spread, so the two are close,
but the ordering on effort favours the classical method.

What survives is narrower and sturdier than what was first written: \emph{at an
appropriate effort weight, a learned policy achieves load-sharing fidelity
comparable to a properly tuned conventional controller, while reaching
substantially more accurately.}

\section{The speed argument, withdrawn}

The original motivation was inference cost. First measurement: classical
\SI{2165}{\micro\second} against the network's \SI{295}{\micro\second}, a
$7.3\times$ advantage.

\begin{trap}
Two objections, and together they are fatal.

\textbf{It timed the wrong computation.} The \SI{2165}{\micro\second} is the
Newton--Raphson solve of Hill fibre--tendon equilibrium, which is part of
\emph{forward simulation}. The question this thesis asks is answered by
\emph{load-sharing optimisation} --- a different calculation entirely. The two had
been conflated.

\textbf{Both sides were unoptimised Python.} Re-timed like-for-like, with
moment-arm extraction charged to the classical side:

\begin{center}
\begin{tabular}{@{} L{7.6cm} R{2.4cm} R{2.4cm} @{}}
\toprule
& \textbf{mean (\si{\micro\second})} & \textbf{p99} \\
\midrule
Network forward pass & 288.8 & 499.9 \\
Classical, general solver (SLSQP) & 2370.5 & 5133.1 \\
\textbf{Classical, direct solve $+$ moment arms} & \textbf{23.9} & --- \\
\bottomrule
\end{tabular}
\end{center}

Against a general-purpose solver the network is $8.2\times$ faster. Against a
direct solve of the actual problem it is \textbf{about twelve times slower}.

\textbf{The latency argument is withdrawn.} The $7.3\times$ reflected the cost of
one Python implementation, not a property of the two approaches.
\end{trap}

\textbf{What survives, honestly stated.} A qualification runs the other way: the
box constraints are active in 100\,\% of sampled states, so the closed-form solve
is not exactly valid and a constrained active-set solver would cost more than
\SI{23.9}{\micro\second}. It would not cost \SI{2370}{}. The defensible statement
is that the two are \emph{comparable} in cost.

What genuinely survives is the \emph{shape} of the cost, not its magnitude: the
network's runtime is fixed and branch-free regardless of biomechanical state,
where the iterative solver's depends on convergence. Where a hard real-time bound
matters more than mean latency, that retains value. It is a considerably weaker
claim than a speed-up, and it is the one made.

The policy contains 393{,}750 parameters --- \SI{1.5}{\mega\byte} single-precision,
\SI{385}{\kilo\byte} quantised to eight bits. These figures are independent of
policy quality, and specific to the SAC architecture (PPO's is 154{,}131).

\section{Muscle synergies, and the weakest of the three analyses}

Motor-control theory holds that the nervous system manages redundant musculature
through a few co-activated groups rather than muscle by muscle. Non-negative
matrix factorisation recovers that structure; VAF measures how much it explains.

\textbf{Two findings and one caution.}

\textbf{The classical solution is uniformly more low-dimensional.} Static
optimisation needs four or five synergies to reach 90\,\% VAF where the learned
policies need six to nine. This is co-contraction expressed in synergy space:
driving an antagonist independently of its agonist requires an extra dimension
that the classical solution never uses.

\textbf{Agreement between learned and classical synergy structure tracks training
quality} --- 0.676 for default SAC against 0.851 for tuned SAC at 300k.

\begin{trap}
\textbf{But this statistic is unstable, and the first version of the result
exploited that without meaning to.} Varying the episode count for one fixed
policy gave 0.873, 0.805, 0.950, 0.810, 0.739 at 8, 12, 15, 20, 25 episodes --- a
spread of 0.21, and non-monotone, so it is sampling noise, not convergence.

The VAF over the same range moved only 0.818--0.870, which localises the problem:
the instability is in the \emph{basis matching}, not the decomposition. NMF has no
unique solution, and two independently fitted bases inherit that ambiguity.

Re-estimated over six random 15-episode subsets, the reported 0.950 becomes
$0.851 \pm 0.129$ --- and 0.950 sits near the top of its own resampled range
$[0.708, 0.979]$.

\textbf{A separate methodological error in the same statistic.} The similarity was
originally a best match per synergy, which lets several extracted synergies claim
the same reference while others go unmatched. On real data only three of four
references were distinctly matched, inflating the score by roughly 0.10. All
figures now use a one-to-one (Hungarian) assignment, which is what a
correspondence between sets actually means.
\end{trap}

\textbf{Confidence: low, and labelled as such.} The ordering survives --- the gap
of 0.175 exceeds the pooled spread --- so the qualitative claim stands. But this
is the weakest of the three convergent analyses and should be weighted
accordingly: the effort ratio (five seeds, non-overlapping distributions) and the
co-contraction index are far better determined than a statistic whose sampling
spread is comparable to the effect it measures.

The comparison against \emph{human} synergies (0.62--0.80) is not interpreted at
all: the reference loadings are a documented qualitative approximation over this
model's muscle ordering, not digitised from the original publication.

\section{Confidence: what is solid, and what is not}

\begin{center}
\begin{tabular}{@{} L{7.0cm} C{2.4cm} L{4.6cm} @{}}
\toprule
\textbf{Claim} & \textbf{Confidence} & \textbf{Basis} \\
\midrule
$\gamma$ causes the improvement & \textbf{high} & two-sided ablation, 4 alternatives excluded, mechanism predicts timescale \\
Replay ratio was 0.233 & \textbf{high} & directly counted; run-time arithmetic closes \\
Load-sharing pattern agrees & \textbf{high} & 5 seeds, proper null, machine-precision validity check \\
Excess effort is co-contraction at the elbow & \textbf{high} & two independent analyses localise it identically \\
Raising $w_2$ closes the gap & \textbf{high} & predicted in advance, 5 seeds, non-overlapping distributions, variance collapses \\
Learned vs conventional trade-off & \textbf{moderate} & both sides tuned, but $n = 1$ controller at one gain pair \\
Algorithm ordering & \textbf{moderate} & 3 seeds, no significance test, 100k snapshot, untuned \\
Synergy agreement tracks quality & \textbf{low} & sampling spread comparable to the effect \\
Anything about \emph{human} motor control & \textbf{none} & no EMG, no subject, no measured force \\
\bottomrule
\end{tabular}
\end{center}

\subsection{The five limitations that matter most}

\textbf{1. No ground truth exists anywhere in this work.} Both sides are computed
in the same simulator on the same muscle model. Where they agree, two
computations over one approximation agree. And static optimisation is itself a
\emph{theory} of load sharing --- well established, still contested. A controller
that disagrees with it is not thereby shown to be unphysiological.

\textbf{2. One task.} A single reaching movement in one workspace. Nothing here
establishes that any finding transfers to another task --- including the $\gamma$
result, whose mechanism explicitly depends on this task's duration.

\textbf{3. Three seeds, no significance testing.} Where five seeds exist they are
used; where three do, no test is claimed, because three cannot support one.

\textbf{4. ``Complete Hill model'' overstates what is active.} The parallel
elastic element returns exactly zero in every sampled state, and the
force--length curve is exercised only on its ascending limb.

\textbf{5. Trajectory realism is not assessed.} The comparison validates load
sharing along a given path, never the path itself.

\section{Questions people actually ask}

\textbf{Did the DRL work?} On reaching: yes --- a hand starting \SI{0.80}{\meter}
away ends within \SI{0.107}{\meter} and holds. On the thesis question: it
reproduces \emph{which} muscles act and in \emph{what proportion} ($r = 0.94$ at
$w_2 = 20$), and it reproduces the \emph{magnitude} too once the objective asks
for it.

\textbf{Is \SI{0.107}{\meter} good?} It is not human-level, and this document does
not claim it is. For context: the untuned baseline was \SI{0.26}{\meter}, a
properly tuned conventional controller reaches \SI{0.165}{\meter}, and a
controller given \emph{full privileged state} has a median final error of
\SI{0.08}{\meter}. The learned policy, acting only on proprioception, sits between
the conventional controller and the privileged one.

\textbf{Why is the success rate so low?} Because it demands \SI{2}{\centi\meter}
\emph{and} settled, and 4\,\% of fifty episodes is two episodes. Across five
identical runs it read 0, 4, 6, 8 and 20\,\%. It is noise, and it is reported as
noise.

\textbf{Should we use DRL instead of classical calculation?} Not for this, on this
evidence. It is not faster (that claim is withdrawn), and it needs the effort
objective supplied explicitly, which is the classical criterion --- so it does not
remove the modelling assumption, it relocates it. What it offers is a fixed,
branch-free inference cost and better endpoint accuracy than a conventional
controller. That is a trade-off, not a replacement.

\textbf{What is the single most useful thing here?} That the discount factor,
matched to the timescale of the behaviour, mattered about three times more than
the choice of algorithm --- and that this was invisible until an ablation was run,
having been confidently misattributed to something else first.

\textbf{Why does the document keep describing its own mistakes?} Because every
one of them changed a number that had already been written down, and in every
case the uncorrected number was the more flattering one. A result that has
survived a serious attempt to break it is worth more than one that has not been
tested; showing the attempts is how a reader can tell which kind they are
holding.
